\begin{pastexam}[label={px:ntut-113-prog-01}]{北科大資工所 113 學年度 第 1 題}
（程式填空）

\textbf{Problem 1 [10\%] [each 2\%]}

Suppose that the outputs of the following C program are as follows:\\
\hspace*{1em}ID: 112590011, GPA: 3.95\\
\hspace*{1em}ID: 112590025, GPA: 3.80\\
\hspace*{1em}ID: 112590001, GPA: 3.75

Please trace the program and fill the blanks problems 1-1$\sim$1-5with correct statements.
\begin{minted}[linenos, numbersep=1.5mm, xleftmargin=2em, fontsize=\small, breaklines]{c}
#include <stdio.h>
#include <stdlib.h>
typedef struct Student {
      int studentID;
      double GPA;
      ______________;                               /*Problem1-1*/
} Student_t;
void printStudentInfo(Student_t *current) {
      while (current != NULL) {
            printf(" ID: %d, GPA: %.2f \n", current->studentID, current->GPA);
            current = current->next;
      }
}
Student_t* addStudent(Student_t* head, int studentID, double GPA) {
      Student_t* newStudent = (Student_t*)malloc(sizeof(Student_t));
      if (______________ ) {                        /*Problem1-2*/
            printf("Memory allocation failed.\n");
            exit(EXIT_FAILURE);
      }
      newStudent->studentID = studentID;
      newStudent->GPA = GPA;
      ________________;                             /*Problem 1-3 */
      return newStudent;
}
void freeList(Student_t* head) {
      Student_t* current;
      while (head != NULL) {
            ____________;                           /*Problem 1-4 */
            free(current);
            head = head->next;
      }
}
void main() {
      ________________;                             /*Problem 1-5 */
      head = addStudent(head, 112590001, 3.75);
      head = addStudent(head, 112590025, 3.80);
      head = addStudent(head, 112590011, 3.95);
      printStudentInfo(head);
      freeList(head);               // Free memory
}
\end{minted}
\end{pastexam}

\begin{pastsoln}
答案：
\begin{itemize}
\item 1-1：\texttt{struct Student *next}
\item 1-2：\texttt{newStudent == NULL}
\item 1-3：\texttt{newStudent->next = head}
\item 1-4：\texttt{current = head}
\item 1-5：\texttt{Student\_t *head = NULL}
\end{itemize}

\textbf{1-1}。 第 11 行用到 \texttt{current->next}，結構裡要有指向下一個節點的指標。 這時 \texttt{typedef}
還沒結束，\texttt{Student\_t} 這個名字還不能用，要寫成 \texttt{struct Student *next}（資結書第~5.1~節的
自我參考結構）。

\textbf{1-2}。 \texttt{malloc} 失敗時回傳 \texttt{NULL}，條件是 \texttt{newStudent == NULL}（或
\texttt{!newStudent}）。

\textbf{1-3、1-5}。 \texttt{main} 依序加入 $112590001$、$112590025$、$112590011$，輸出的順序剛好相反，
最後加入的在最前面，所以 \texttt{addStudent} 是把新節點插在串列\textbf{開頭}：新節點的 \texttt{next}
指向原本的開頭，再回傳新節點當作新的開頭（資結書虛擬碼~5.2 的頭端插入），$O(1)$。 \texttt{main} 第 35 行
起直接使用 \texttt{head}，第 34 行要宣告它並設成空串列 \texttt{NULL}；第一次插入時新節點的 \texttt{next}
就是 \texttt{NULL}，成為串列的結尾。 三次插入後：
\[
\text{head} \to 112590011 \to 112590025 \to 112590001 \to \text{NULL},
\]
\texttt{printStudentInfo} 從頭印到尾，得到題目的輸出（\texttt{printf} 的格式字串開頭有一個空白，題目列出
的輸出省略了）。

\textbf{1-4}。 第 29 行 \texttt{free(current)}，\texttt{current} 要先指向這一輪要釋放的節點，也就是
\texttt{head}。

\textbf{原程式的錯誤}。 照這樣填，第 29、30 行先 \texttt{free(current)}（就是 \texttt{head} 指的節點），
再讀 \texttt{head->next}，讀的是已經釋放的記憶體（use after free），行為未定義；多數情況下剛好還讀得到
舊值，看不出問題，但用 AddressSanitizer 檢查會直接報錯。 正確的寫法是先往前走再釋放：
\begin{minted}[fontsize=\small]{c}
while (head != NULL) {
      current = head;
      head = head->next;
      free(current);
}
\end{minted}
第 28 行只能填一個敘述，填不出正確的版本，作答時寫 \texttt{current = head} 並註明這個問題即可。 另外，
\texttt{void main()} 不是標準 C 的寫法，應該寫 \texttt{int main(void)}。
\end{pastsoln}
